\originalsection{附录 A: 扩展讨论}\paperlabel{A}{chive-app:extended_discussion}
\subsection*{A.1 评价代理论证}\paperlabel{A.1}{chive-app:proxy}
任何能够对目标采样的人, 都可以直接运行反事实来回答我们的评价, 所以在此表现好本身没有直接用途. 我们认为, 行为解释方法的评价大体无法避免这一点.

最需要可解释性工具的地方, 是没有其他方法能检查答案的地方. 例如 Sonnet 4.5 和 Mythos 5 的系统卡, 用可解释性工具研究没有可用真值的问题, 如未言说的评价意识、模型对评分者的推理、未言说的负面情绪\cite{chive-anthropic2025sonnet45card,chive-anthropic2026mythos5card}. 没有已知真值, 就难判断该多依赖工具, 大多数例子对工具能信任到何种程度都留有保留.

因此我们认为值得构建代理评价: 有明确真值, 问题形式与现实用途相同（给定异常行为, 解释为何发生）, 同时接受评价表现本身没有直接用途. 在代理上有改善的工具可能意味着可靠性提高, 从而增加对没有东西可检查工具输出的应用案例的信任. 在这里失败, 对难以找出原因的更困难案例是坏信号.

主要保留是: 代理好坏取决于可由提示反事实检查的因素, 与实际使用可解释性工具的因素之间的重合. 一个乐观理由是, 许多反事实归结到模型内部对输入的表征, 如如何理解歧义句, 所以评价仍与内部表征有一些联系.

\subsection*{A.2 假设生成与假设辨别}\paperlabel{A.2}{chive-app:hypothesis_discrimination}
常见论点是, 可解释性工具适合生成假设\cite{chive-frasertaliente2026nla}, 因为输出不足够可靠, 不能完全信任. 这似乎意味着评价不应聚焦假设辨别, 即识别哪个假设正确的能力. Fraser-Taliente 等论文\cite{chive-frasertaliente2026nla}是少数公开利用可解释性工具生成假设、调试模型行为的案例之一. 然而我们检查其发表的例子后发现, 每例中由工具发现的假设, 都是 Opus 只看行为描述就能猜出的假设.

论文的四个案例及工具发现的原因如下:
\begin{itemize}
\item \textbf{语言切换}: 英文提示得到俄文回答. 原因是格式错误的微调资料, 将译成英语的提示配上仍保留原外语的回答.
\item \textbf{误报工具结果}: 模型无声忽略错误计算器输出, 报出正确答案. 题很简单, 模型阅读时已自己计算, 比起错误工具结果更信任自己的计算.
\item \textbf{诗歌中的提前规划}: 模型补全押韵两行诗. 在第一行末,“提前规划”并已表征第二行结尾的候选词, 朝它写作.
\item \textbf{答案反复摇摆}: 尽管自己的推理得到正确答案, 仍反复陈述错误答案. 训练中记住了带错误标签的题, 记忆答案与上下文计算冲突.
\end{itemize}
事后看, 这些原因都不很惊人, 可能只由行为提出. 为测试 LLM 能否猜到, 我们为各行为写不提原因的描述, 要求列出前五个假设, 每模型采样十次. Opus 4.8 评判者报告第一个符合 NLA 论文所述原因的假设排名. 使用 Opus 4.1 与 4.5, 其训练截止时间早于 NLA 论文, 因而不可能只是记住这篇论文.

表 \ref{chive-tab:hypothesis_generation} 显示结果. 每个模型--行为组合中, 发现的假设至少在 60\% 样本里进入前三, 最常是第一猜测. 其他常见例子如评价意识、对奖励信号的推理, 是审计者进入评价场景时已有的假设. 此时工具验证已有假设, 并非生成新假设. 更一般地, 我们预期很少有意外行为是前沿模型在 10--100 次猜测内不能给出近似正确假设的. 若这判断正确, 假设辨别是值得改善的重要能力.

假设生成在其他工作流仍有价值, 如扫描大量记录标出异常工具生成, 或对微调模型实验体的怪异特性提出假设. 这些是真实用途, 但区别于解释某项特定行为, 后者才是本评价目标.

\begin{table}[ht]\centering
\begin{tabular}{lrrrr}\toprule
 & \multicolumn{2}{c}{Opus 4.1} & \multicolumn{2}{c}{Opus 4.5}\\
行为 & 首位 & 前三 & 首位 & 前三\\\midrule
语言切换 & 9/10 & 10/10 & 3/10 & 6/10\\
误报工具结果 & 10/10 & 10/10 & 10/10 & 10/10\\
诗歌规划 & 2/10 & 8/10 & 7/10 & 10/10\\
答案反复摇摆 & 10/10 & 10/10 & 6/10 & 10/10\\\bottomrule
\end{tabular}
\caption{从行为描述生成假设. 对 Fraser-Taliente 等\cite{chive-frasertaliente2026nla}每个假设生成案例, 每格十个样本中, 工具发现的原因出现在模型第一假设（top-1）或前三假设（top-3）的比例.}\label{chive-tab:hypothesis_generation}
\end{table}
